% GENERATED FILE. Edit canonical lessons, then run npm run generate:books.
% canonical-source-hash: fnv1a64:049a6e99efc4499e
% canonical-lessons: ML-C235-kanankai, ML-C235-uppurri, ML-C235-nattellu, ML-C235-karal, ML-C235-vrkka

\chapter{Wrist, Heel, Spine, Liver, Kidney}
\label{ch:ml-wrist-heel-spine-liver-kidney}
\hlchaptermodality{235}

\begin{chapteropening}
\textbf{By the end of this chapter:} \emph{I can say a wrist, a heel, the spine, the liver and a kidney.}

Complete the last lesson of chapter 235: I can say a wrist, a heel, the spine, the liver and a kidney.
\end{chapteropening}

\section[\texorpdfstring{\ml{കണങ്കൈ} (kaṇaṅkai)}{കണങ്കൈ (kaṇaṅkai)}]{\ml{കണങ്കൈ} (kaṇaṅkai) — a wrist}

\label{lesson:ML-C235-kanankai}



\begin{quote}
Before the new one: say the Malayalam for excitement, then the Malayalam for relief.
\end{quote}

\subsection*{You'll want to know: \ml{കണങ്കൈ}}
\textbf{\ml{കണങ്കൈ}} — \emph{kaṇaṅkai} — \textquotedblleft{}a wrist\textquotedblright{}.

\begin{cousinweb}
\textbf{\ml{കണങ്കൈ}} (\emph{kaṇaṅkai}) is native Dravidian, \textquotedblleft{}the joint of the hand\textquotedblright{}: \emph{kaṇ}, \textquotedblleft{}a joint, a node\textquotedblright{} (DEDR 1160), and \emph{kai}, \textquotedblleft{}hand\textquotedblright{} (DEDR 2023).
\end{cousinweb}

\begin{grammarlens}[title={What you've built}]
One more word for this chapter.
\end{grammarlens}

\subsection*{Your turn}
\begin{itemize}
\raggedright
  \item \emph{Say it:} \emph{kaṇaṅkai}
  \item \emph{Say it:} \emph{kaṇaṅkai}, once more
  \item \emph{Say it:} say \emph{kaṇaṅkai}
  \item \emph{Recall:} say the Malayalam for excitement, then the Malayalam for relief, then say \emph{kaṇaṅkai} again
\end{itemize}

\begin{tcolorbox}[breakable,colback=teal!4,colframe=teal!35!black,title={Before you move on}]
What does \ml{കണങ്കൈ} mean? (\textquotedblleft{}A wrist\textquotedblright{}.) Say it once more.
\end{tcolorbox}
\section[\texorpdfstring{\ml{ഉപ്പൂറ്റി} (uppūṟṟi)}{ഉപ്പൂറ്റി (uppūṟṟi)}]{\ml{ഉപ്പൂറ്റി} (uppūṟṟi) — a heel}

\label{lesson:ML-C235-uppurri}



\begin{quote}
Before the new one: say the Malayalam for relief, then the Malayalam for a wrist.
\end{quote}

\subsection*{You'll want to know: \ml{ഉപ്പൂറ്റി}}
\textbf{\ml{ഉപ്പൂറ്റി}} — \emph{uppūṟṟi} — \textquotedblleft{}a heel\textquotedblright{}.

\begin{cousinweb}
\textbf{\ml{ഉപ്പൂറ്റി}} (\emph{uppūṟṟi}) is of uncertain origin; DEDR, the dictionary of the Dravidian family, has no entry for it.
\end{cousinweb}

\begin{grammarlens}[title={What you've built}]
One more word for this chapter.
\end{grammarlens}

\subsection*{Your turn}
\begin{itemize}
\raggedright
  \item \emph{Say it:} \emph{uppūṟṟi}
  \item \emph{Say it:} \emph{uppūṟṟi}, once more
  \item \emph{Say it:} say \emph{uppūṟṟi}
  \item \emph{Recall:} say the Malayalam for relief, then the Malayalam for a wrist, then say \emph{uppūṟṟi} again
\end{itemize}

\begin{tcolorbox}[breakable,colback=teal!4,colframe=teal!35!black,title={Before you move on}]
What does \ml{ഉപ്പൂറ്റി} mean? (\textquotedblleft{}A heel\textquotedblright{}.) Say it once more.
\end{tcolorbox}
\section[\texorpdfstring{\ml{നട്ടെല്ല്} (naṭṭellŭ)}{നട്ടെല്ല് (naṭṭellŭ)}]{\ml{നട്ടെല്ല്} (naṭṭellŭ) — the spine, the backbone}

\label{lesson:ML-C235-nattellu}



\begin{quote}
Before the new one: say the Malayalam for a wrist, then the Malayalam for a heel.
\end{quote}

\subsection*{You'll want to know: \ml{നട്ടെല്ല്}}
\textbf{\ml{നട്ടെല്ല്}} — \emph{naṭṭellŭ} — \textquotedblleft{}the spine, the backbone\textquotedblright{}.

\begin{cousinweb}
\textbf{\ml{നട്ടെല്ല്}} (\emph{naṭṭellŭ}) is native Dravidian, \textquotedblleft{}the middle bone\textquotedblright{}: \emph{naṭu}, \textquotedblleft{}middle\textquotedblright{} (DEDR 3584, with Tamil \emph{naṭu}), and \emph{ellŭ}, \textquotedblleft{}bone\textquotedblright{} (DEDR 839). DEDR 839 compares Sanskrit \emph{eḍūka}.
\end{cousinweb}

\begin{grammarlens}[title={What you've built}]
One more word for this chapter.
\end{grammarlens}

\subsection*{Your turn}
\begin{itemize}
\raggedright
  \item \emph{Say it:} \emph{naṭṭellŭ}
  \item \emph{Say it:} \emph{naṭṭellŭ}, once more
  \item \emph{Say it:} say \emph{naṭṭellŭ}
  \item \emph{Recall:} say the Malayalam for a wrist, then the Malayalam for a heel, then say \emph{naṭṭellŭ} again
\end{itemize}

\begin{tcolorbox}[breakable,colback=teal!4,colframe=teal!35!black,title={Before you move on}]
What does \ml{നട്ടെല്ല്} mean? (\textquotedblleft{}The spine, the backbone\textquotedblright{}.) Say it once more.
\end{tcolorbox}
\section[\texorpdfstring{\ml{കരൾ} (karaḷ)}{കരൾ (karaḷ)}]{\ml{കരൾ} (karaḷ) — the liver}

\label{lesson:ML-C235-karal}



\begin{quote}
Before the new one: say the Malayalam for a heel, then the Malayalam for the spine.
\end{quote}

\subsection*{You'll want to know: \ml{കരൾ}}
\textbf{\ml{കരൾ}} — \emph{karaḷ} — \textquotedblleft{}the liver\textquotedblright{}.

\begin{cousinweb}
\textbf{\ml{കരൾ}} (\emph{karaḷ}) is native Dravidian, one word with Kannada \emph{karuḷ} (DEDR 1274). Its older sense covers the heart and the inner organs together.
\end{cousinweb}

\begin{grammarlens}[title={What you've built}]
One more word for this chapter.
\end{grammarlens}

\subsection*{Your turn}
\begin{itemize}
\raggedright
  \item \emph{Say it:} \emph{karaḷ}
  \item \emph{Say it:} \emph{karaḷ}, once more
  \item \emph{Say it:} say \emph{karaḷ}
  \item \emph{Recall:} say the Malayalam for a heel, then the Malayalam for the spine, then say \emph{karaḷ} again
\end{itemize}

\begin{tcolorbox}[breakable,colback=teal!4,colframe=teal!35!black,title={Before you move on}]
What does \ml{കരൾ} mean? (\textquotedblleft{}The liver\textquotedblright{}.) Say it once more.
\end{tcolorbox}
\section[\texorpdfstring{\ml{വൃക്ക} (vṛkka)}{വൃക്ക (vṛkka)}]{\ml{വൃക്ക} (vṛkka) — a kidney}

\label{lesson:ML-C235-vrkka}



\begin{quote}
Before the new one: say the Malayalam for the spine, then the Malayalam for the liver.
\end{quote}

\subsection*{You'll want to know: \ml{വൃക്ക}}
\textbf{\ml{വൃക്ക}} — \emph{vṛkka} — \textquotedblleft{}a kidney\textquotedblright{}.

\begin{cousinweb}
\textbf{\ml{വൃക്ക}} (\emph{vṛkka}) $\leftarrow$ Sanskrit \emph{vṛkka}, \textquotedblleft{}kidneys\textquotedblright{}, taken over whole as a tatsama (CDIAL 12064).
\end{cousinweb}

\begin{grammarlens}[title={What you've built}]
One more word for this chapter.
\end{grammarlens}

\subsection*{Your turn}
\begin{itemize}
\raggedright
  \item \emph{Say it:} \emph{vṛkka}
  \item \emph{Say it:} \emph{vṛkka}, once more
  \item \emph{Say it:} say \emph{vṛkka}
  \item \emph{Recall:} say the Malayalam for the spine, then the Malayalam for the liver, then say \emph{vṛkka} again
\end{itemize}

\begin{tcolorbox}[breakable,colback=teal!4,colframe=teal!35!black,title={Before you move on}]
What does \ml{വൃക്ക} mean? (\textquotedblleft{}A kidney\textquotedblright{}.) Say it once more.
\end{tcolorbox}
